% Section 1.5: 论文结构安排

本文共分六章，整体沿「背景与问题 $\to$ 理论基础 $\to$ 数据构建 $\to$ 方法设计 $\to$ 实验验证 $\to$ 总结展望」的主线展开，各章内容要点如下：

\begin{itemize}
    \item \textbf{第一章\quad 引言}：阐述研究背景与意义，综述国内外现状，提炼稳定性、规划能力、高效适配与训推协同四大关键挑战，并给出研究目标、主要贡献与全文结构。
    \item \textbf{第二章\quad 相关理论与技术基础}：介绍大语言模型与代码生成、参数高效微调（重点为 LoRA）、零样本 / 少样本思维链推理，以及可执行正确性与 Pass@k 等评测指标。
    \item \textbf{第三章\quad 数据集构建与预处理}：说明 KodCode 数据来源与任务形式化，给出格式规范化、执行一致性校验、思维步骤与算法标签标注等完整流水线，并给出训练 / 评测集划分。
    \item \textbf{第四章\quad 模型方法与实验框架设计}：给出训练侧 $\times$ 推理侧的二因素析因框架，定义 $H_1$--$H_{10}$ 完整实验矩阵，并逐一说明 Base、LoRA 微调、零样本 / 少样本 CoT 的提示构造与训练目标。
    \item \textbf{第五章\quad 实验设计与结果分析}：在统一评测协议下报告十组对照的 Pass@k 指标，估计训练侧与推理侧的主效应、交互效应与少样本边际效应，并结合典型错例开展归因分析。
    \item \textbf{第六章\quad 总结与展望}：总结全文研究内容与核心结论，梳理创新点与局限性，并对算法代码生成的后续研究方向进行展望。
\end{itemize}